\section{考察・限界・倫理}
\label{sec:discussion}

\subsection{注釈費用の移動先}
\label{sec:amortization}

本手法は幾何の入手を無償にはしない。
人手の費用を\emph{シーン単位へ移す}。
従来の作業は画像ごとに点を打つか校正する。
本手法では、動画の取得、再構成、
そしてシーンごとに一度のアライメント検証を支払い、
その後の視点配置、姿勢、内部パラメータ、点の同一性、
線と領域の教師、深度、幾何妥当性がすべて自動で付く。
節約は生成した視点数と多様性に比例して増える。
斜め、高所、軸外、部分的なコート、
複数コートといった、局所的な手掛かりでは解けない視点を
安価に追加できることがこの設計の要点である。

\subsection{3DGS は必要条件だが十分条件ではない}
\label{sec:gs_limit}

新規視点合成だけではメートルも意味も得られない。
綺麗だが整合していないガウスシーンは、
並進や物理点の同一性を監督できない。
逆にコート変換だけがあっても、
新しい姿勢での写実的な画素は得られない。
寄与はこの二つを一つの変換グラフで結び、
尺度、トポロジー、holdout 証拠が弱いときに
その結合を棄却する点にある。

また、レンダ結果は撮影会場に強く固定される。
浮遊物、ぼけ、一時的な物体の欠落、
撮影範囲の外側での外挿の不安定さは残る。
軌道を復元カメラのオフセットに錨付けし、
凸包で解析的に制限するのはこのためである。
それでも、凸包制約は画質の保証ではない。
未観測領域を跨ぐ経路や、
高さ方向の外挿が残りうる。

\subsection{アライメントの誤差源}
\label{sec:alignment_failure}

アライメントは学習済みの線検出器を入力に取る。
したがって線検出器の誤りがそのまま誤差源になる。
疎な再構成点が地面を十分に支えない場合、
反射や遮蔽で線が欠ける場合、
隣接する似たコートが 90 度やコート同一性の誤りを誘発する場合がある。
共通尺度、全テンプレート、トポロジー、holdout の各 gate は
この伝播を抑えるが、
受理されたすべての変換が厳密に正しいことを証明するものではない。

\subsection{カメラモデルの範囲}
\label{sec:camera_model}

本稿の検出器は焦点距離の等方性を仮定し、
主点を補正後の対象から受け取る。
歪み、skew、主点の移動は予測しない。
未対応の変換は明示的に棄却されるため
教師の一意性は保たれるが、
広いカメラモデルにはヘッドと投影損失の拡張が要る。
カメラ端の正準化も
\(I\) と \(R_z(\pi)\) の二択を仮定しており、
中央面近傍の強く曖昧な視点は
割り当てずに除外しなければならない。

\subsection{可視性と教師の意味}
\label{sec:visibility_semantics}

投影したすべてのコート構造を予測させるのか、
画像上で見えるマーキングだけを予測させるのかは、
先行研究の可視領域分割と比較可能でなければならない。
本稿はカメラ前方、画面内、描画支持、教師有効の四状態を保存し、
教師有効の定義をデータ側で固定する。
しかし、遮蔽された物理線をどう扱うかの選択は
実装の設計判断であり、
異なる選択をした研究との数値比較には注意を要する。

\subsection{主張の及ぶ範囲}
\label{sec:claim_boundary}

本稿の実験的な主張は、
実写 held-out validation と合成 trajectory-group test で
交差モデル比較を実測したこと、
および視点拡張仮説を切り分ける比較設計までである。
合成 test の差は大きいが、
比較した二つのモデルはアーキテクチャも学習データも異なる。
したがって差を視点拡張の因果効果として読むことはしない。
既知会場の未学習視点と、
未知会場への汎化は別の設定である。
軌道グループ分離は前者の一部を支えるが、
後者を自動的には保証しない。
最終の評価画像を再構成やカメラ生成の入力に使う場合、
その事実を開示し、
厳密な未見画像・未見会場の評価と区別する。
検証用の実写データをカメラ生成の調整に使うなら、
その集合を最終 test から分離する。

実写のみの同一アーキテクチャ、
同枚数の狭域・広域、既知と未知の会場、
複数 seed が揃うまで、
結論は仮説の水準にとどまる。

\paragraph{実写で差が小さいことをどう読むか}
実写 held-out validation では \Method が僅かに上回るが、
差は PCK@0.01 で \resRealOursPckOne 対 \resRealTcdPckOne である。
放送主カメラの分布に近い条件では、
公開実装ベースラインも十分に機能する。
本研究の主張は、
その分布の外側を安価に増やせるという点にあり、
既存手法を分布内で置き換えることではない。
実写側の差を根拠に「より優れた検出器」と結論しない。

\paragraph{古典的手法の位置づけ}
\cref{sec:results_classical} の補助 probe は、
固定版を OpenCV 4 で動かしたときの挙動を記録したものである。
標本は主比較と別であり、
また \texttt{B00} の test は二面のコートが写るシーンであるため、
単一コート前提の実装が不利な条件を含む。
これは「古典的手法が合成データ全般に弱い」という因果主張ではない。
multi-court と非放送視点に対する補助的な robustness の証拠として提示する。

\subsection{計算量}
\label{sec:compute}

再構成と新規視点レンダリングは学習の前に費用が発生する。
一つのシーンから多数の視点と複数の下流タスクを生成できるため、
この費用は償却されるが、
シーン数が少ない場合や
会場ごとに再構成し直す場合は割高になる。
本稿は計算量の詳細な比較を実測していないため、
数値としての計算優位は主張しない。

\subsection{倫理とデータ管理}
\label{sec:ethics}

ソース動画と再構成された外観には、
人物や会場を特定しうる情報が含まれうる。
公開実装のソフトウェアライセンスと、
映像・再構成物の公開条件は別である。
データセットやレンダ結果を公開する場合は、
映像の権利とプライバシーの制約を
ソフトウェアのライセンスとは独立に確認する。
本研究が使う公開実装ベースラインには
再配布を許す条件が確認できないため、
そのコードや重みを再配布しない。
実験に用いる重みとデータは、
本文および付録に記録した同定情報の範囲で参照する。
